\documentclass[12pt]{article}
\usepackage[utf8]{inputenc}
\usepackage[T1]{fontenc}
\usepackage{hyperref}
\title{Parawen Language Resources for K-12 Curriculum Development}
\author{Compiled by Claude Sonnet 5.5 (claude-sonnet-5-5) for Andr\'as Kornai}
\date{}
\begin{document}
\maketitle

\section{Introduction}
Parawen (ISO 639-3 \texttt{prw}) is a language of Papua New Guinea with about 430 native speakers (a 1981 figure cited by Wikipedia); the 2000 census counted 760 ethnic members. Sources disagree on its location (Madang vs.\ Morobe Province) and on its classification (Wikipedia: Papuan; Joshua Project: Austronesian). Wikipedia notes a small phonemic inventory (nine consonants, three vowels), tracing to Andrew Pick's 2020 dissertation, which I did not locate directly. Almost nothing educational was found.

\section{Primary Educational Resources}
No primers, textbooks or school programmes were found.

\section{Linguistic and Reference Materials}
\begin{itemize}
\item Wikipedia, Parawen language: \url{https://en.wikipedia.org/wiki/Parawen_language}
\item Trans New Guinea database, source Zgraggen 1980 (Madang languages; page appeared in results, contents not verified): \url{https://transnewguinea.org/source/zgraggen-1980-nadelbert?page=120}
\item Endangered Languages Project bibliography entry: \url{https://www.endangeredlanguages.com/lang/10740/bibliography?hl=zh-cn}
\end{itemize}

\section{Religious and Community Resources}
\begin{itemize}
\item Joshua Project people profile (Scripture status: translation started): \url{https://joshuaproject.net/people_groups/14318}
\item Joshua Project language page: \url{https://joshuaproject.net/languages/prw}
\end{itemize}

\section{Digital and Audio Resources}
None found.

\section{Curriculum Development Context}
Published material is scarce, so curriculum work would depend on new documentation with the community.

\end{document}